\section{Conclusão}

O trabalho caracterizou um sistema LIT discreto a partir de conjuntos de dados de
entrada e saída e chegou a um modelo capaz de descrever bem o seu comportamento na
faixa de frequência excitada.

No caso de entrada branca, a análise no tempo e na frequência indicou um sistema do
tipo passa-baixas. A correlação cruzada forneceu uma estimativa direta da resposta ao
impulso, e o modelo FIR identificado por mínimos quadrados alcançou bom ajuste tanto
no treino quanto na validação, com resíduos próximos de um ruído branco. A coerência
mostrou em quais faixas a relação entrada--saída pode ser considerada confiável.

No caso de entrada colorida, a entrada concentrou energia nas frequências baixas e
deixou as altas pouco excitadas. A comparação com o caso branco mostrou que a
identificação piora justamente nessas faixas, e que a coerência antecipa essa perda de
confiabilidade. Ficou claro também que um bom valor de ajuste global não garante um
modelo correto em todas as frequências.

No caso do conjunto de teste, a aplicação do modelo permitiu calcular os resíduos e
localizar uma mudança de regime, além de um aumento da variância e de impulsos
isolados. A comparação mostrou que a anomalia se manifesta com mais clareza nos
resíduos do que no sinal de saída, o que reforça o uso dos resíduos como ferramenta de
monitoramento.

De modo geral, a autocorrelação, a densidade espectral de potência, o espectro
cruzado, a coerência e os resíduos atuaram de forma complementar. As três primeiras
descrevem o sinal e a relação entrada--saída, a coerência indica onde essa relação é
confiável, e os resíduos verificam o quanto resta por explicar depois de aplicado o
modelo. Esse conjunto de ferramentas permitiu avaliar a qualidade e os limites do
modelo identificado.

\section*{Reprodução dos resultados}
\addcontentsline{toc}{section}{Reprodução dos resultados}

Os resultados podem ser reproduzidos a partir do repositório do trabalho. A análise
completa está no caderno \texttt{notebooks/resolucao\_trabalho.ipynb}. As figuras e os
valores numéricos deste relatório são gerados pelo script
\texttt{relatorio/gerar\_figuras.py}, que lê os mesmos arquivos de dados e grava as
figuras na pasta \texttt{relatorio/figuras}. As dependências estão listadas em
\texttt{requirements.txt}.
